% Erwin Poussi — CV
% Single-column LaTeX CV matching the previous resume.pdf layout.
\documentclass[10pt,letterpaper]{article}

\usepackage[top=0.4in,bottom=0.4in,left=0.55in,right=0.55in]{geometry}
\usepackage[T1]{fontenc}
\usepackage[utf8]{inputenc}
\usepackage{lmodern}
\usepackage{microtype}
\usepackage{enumitem}
\usepackage{titlesec}
\usepackage{hyperref}
\usepackage{xcolor}
\usepackage{tabularx}
\usepackage{parskip}

\hypersetup{
  colorlinks=true,
  urlcolor=blue!60!black,
  linkcolor=blue!60!black,
  pdftitle={Erwin Poussi - Resume},
  pdfauthor={Erwin Poussi}
}

\pagestyle{empty}
\setlength{\parindent}{0pt}

% Section header: ALLCAPS, bold, with rule below
\titleformat{\section}
  {\normalsize\bfseries\uppercase}
  {}{0em}{}[\vspace{-3pt}\titlerule\vspace{1pt}]
\titlespacing*{\section}{0pt}{4pt}{1pt}

% Bullet list tuning
\setlist[itemize]{leftmargin=1.1em,itemsep=0pt,topsep=1pt,parsep=0pt}
\setlength{\parskip}{1pt}

% Entry helper: \entry{title}{location}{date}
\newcommand{\entry}[3]{%
  \vspace{0pt}\noindent\textbf{#1} \textit{#2}\hfill #3\par}

% Sub-entry (research project): \subentry{title}{date}
\newcommand{\subentry}[2]{%
  \vspace{0pt}\noindent\textbf{#1}\hfill #2\par}

\begin{document}

%---------- Header ----------
\begin{center}
  {\LARGE \textbf{Erwin POUSSI}}\\[2pt]
  \href{mailto:erwinpi@stanford.edu}{erwinpi@stanford.edu} \ \ $\vert$\ \ \textbf{US:} +1 650 250 7491 \ \ $\vert$\ \ \textbf{FR:} +33 (0)7 49 62 90 59 \ \ $\vert$\ \ \href{https://rwin2.github.io/erwinpoussi.github.io/}{Website} \ \ $\vert$\ \ \href{https://github.com/Rwin2}{GitHub} \ \ $\vert$\ \ \href{https://www.linkedin.com/in/erwin-poussi-a05099215/}{LinkedIn}\\[2pt]
  \textit{MS in Aero/Astro at Stanford --- robot learning, motion planning, LLM fine-tuning, sim-to-real transfer.}
\end{center}

%---------- Education ----------
\section{Education}

\entry{Stanford University,}{Palo Alto (CA)}{2025 -- 2027}
\begin{itemize}
  \item MS in Aeronautics \& Astronautics \hfill GPA: \textbf{4.1/4.0}
  \item Relevant coursework: \textit{Principles of Robot Autonomy (AA274A/CS237A), Introduction to Robotics (CS223A), Decision Making under Uncertainty (AA228/CS238), Optimal and Learning-Based Control (AA203), Deep Learning \& Reinforcement Learning (CS224R).}
\end{itemize}

\entry{Ecole Polytechnique,}{Palaiseau (FRA)}{2020 -- 2025}
\begin{itemize}
  \item Master's level in Mechanical Engineering, concentration in Autonomy, minor in Entrepreneurship. GPA: \textbf{3.96/4.0}
  \item Relevant coursework: \textit{Advanced Machine Learning for Autonomous Agents, Optimization and Control, Statistics.}
\end{itemize}

%---------- Work ----------
\section{Work Experience}

\entry{Stanford Multi-Robot Systems Lab}{Palo Alto, CA}{Jan.\ 2025 -- Present}
\begin{itemize}
  \item \textbf{Foundation Model Training --- Imitation Learning \& Sim-to-Real Transfer:} Trained vision-language-action (VLA) policies for autonomous drone navigation via behavioral cloning and DAgger.
  \item Implemented RRT* trajectory planning + MPC expert policy for data generation; grounded language instructions to 3D Gaussian Splat scenes via CLIPSeg; trained BC + DAgger policies, reaching \textbf{90\% navigation success} (up from 52\%). \textit{(\href{https://rwin2.github.io/erwinpoussi.github.io/project/language-steered-drones/}{project page})}
  \item \textbf{GAIT --- Agentic Scene Understanding for Open-Vocabulary Navigation} (Jun.--Sep.\ 2026): Built a graph-memory architecture for VLA models leveraging online semantic scene graph construction to enable open-vocabulary navigation with interrupt-triggered re-planning. \textit{(\href{https://rwin2.github.io/erwinpoussi.github.io/project/gait/}{project page})}
\end{itemize}

\entry{Stanford School of Medicine, Qiu Lab}{Palo Alto, CA}{Oct.\ 2025 -- Present}
\begin{itemize}
  \item \textbf{ML Researcher --- \href{https://www.pantheon-os.com/}{Pantheon-CLI}:} Multi-agent LLM framework for autonomous scientific discovery with tool-augmented inference for genomic analysis.
  \item Implemented actor-critic RL for gene panel selection in scRNA-seq; fine-tuned Qwen3 LLMs via SFT and RLOO to power a self-evolving code-generation agent backbone. \textit{(\href{https://rwin2.github.io/erwinpoussi.github.io/publication/xu-2026-pantheonos/}{project page})}
  \item \textbf{Virtual Embryo Challenge --- 3D Convolutional Models for Embryogenesis} (Jun.--Sep.\ 2026): Training foundational 3D convolutional models to reproduce embryogenesis, predicting cell-level gene expression dynamics and morphogenesis; co-building the \href{https://virtualembryo.ai/challenge}{Virtual Embryo Challenge}. \textit{(\href{https://rwin2.github.io/erwinpoussi.github.io/project/virtual-embryo/}{project page})}
\end{itemize}

\entry{Stevens Institute of Technology, RRG}{Hoboken, NJ}{Apr.\ 2025 -- Aug.\ 2025}
\begin{itemize}
  \item \textbf{Research Intern --- ML-Based Aerodynamic Modeling:} Modeled parachute broadcloth permeability under rarefied conditions using a multi-fidelity ML framework combining random forests and genetic algorithms for parameter inference from experimental data; achieved \textbf{$<$1\% mean error}. \textit{(\href{https://rwin2.github.io/erwinpoussi.github.io/project/parachute-opt/}{project page})}
\end{itemize}

%---------- Research ----------
\section{Selected Research}

\subentry{Reinforcement Learning for Algorithmic Trading}{Jan.\ 2025 -- Apr.\ 2025}
\begin{itemize}
  \item Built a custom high-frequency market simulator and trained agents with PPO, Actor-Critic, and DQL, focusing on reward-free exploration, robustness to distribution shifts, and real-time decision-making.
\end{itemize}

\subentry{Guidance System for Autonomous Visual-Based Docking --- Nyx Spacecraft (CSEP)}{Sep.\ 2023 -- May 2024}
\begin{itemize}
  \item Undergraduate thesis on vision-based navigation for the Nyx spacecraft (selected by Axiom for ISS support). Supervised by The Exploration Company with Polytechnique Space Center. \textit{(\href{https://rwin2.github.io/erwinpoussi.github.io/project/visual-based-docking/}{project page})}
\end{itemize}

%---------- Skills ----------
\section{Technical Skills}

\noindent\textbf{Programming \& Frameworks}: Python, PyTorch, CUDA, ROS2, MuJoCo, Hugging Face / Transformers, ACADOS, GSplat, GTSAM, CLIPSeg, Habitat-Sim, W\&B\\
\textbf{Methods}: Behavioral cloning, DAgger, PPO, Actor-Critic, DQL, RRT*, MPC, SFT, RLOO, RLVR, LLM fine-tuning, imitation learning, sim-to-real transfer, motion planning, trajectory optimization\\
\textbf{Languages}: French (native), English (TOEFL 109), Mooré (fluent), German (basic)

%---------- Publications ----------
\section{Publications}

\noindent\textbf{Xu, W., Poussi, E.,} et al. \textit{PantheonOS: An Evolvable Multi-Agent Framework for Automatic Genomics Discovery.} Preprint, 2026. \href{https://doi.org/10.64898/2026.02.26.707870}{DOI: 10.64898/2026.02.26.707870}

%---------- Extracurriculars ----------
\section{Extracurricular Activities}

\noindent\textbf{Volunteering:} Stanford BGSA Fundraising Chair (2025--Present) $\cdot$ Cheer Up cancer nonprofit (2023--2025) $\cdot$ XProjets consulting PM (2024--2025) $\cdot$ Math/physics tutor (2022--2023)\\
\textbf{Sports:} Handball (France's intercollegiate championship), Track \& Field, Weightlifting \quad \textbf{Arts:} Drawing, Painting, Dancing

\noindent\textbf{Awards:} French Government Excellence Scholarship (Ecole Polytechnique) $\cdot$ 2\textsuperscript{nd} Place, 2021 Youth Olympiads, French Legion of Honor $\cdot$ 1\textsuperscript{st} nationwide, 2020 Burkina Faso Scientific A-levels $\cdot$ 2020 Model Student Award $\cdot$ 2020--2018 Excellence Awards

\end{document}
